\clearpage
\fancyhead[L]{\small MÜLLER. BERECHNUNG DER SONNENFINSTERNISSE.}
\hypertarget{native-596-lower}{}
\begin{center}\small Göttingische gelehrte Anzeigen. Stück 20. Seite 200. 3. Februar 1816.\\[3pt]\rule{.72\linewidth}{.3pt}\end{center}

\emph{Ausführliche Anleitung zur trigonometrischen Berechnung der an einem gegebenen Ort der Erdoberfläche sichtbaren Sonnenfinsternisse}, nach zwei verschiedenen, sehr genauen Methoden, erläutert anhand der Bestimmung der Erscheinungen der großen Sonnenfinsternis vom 19. November 1816 für den Nürnberger Meridian, von \textsc{Johann Wolfgang Müller}, Professor der Mathematik am Königlichen Gymnasium in Nürnberg. Sulzbach. In der Kunst- und Buchhandlung des Kommerzienrats J. E. \textsc{Seidel}. 128 Seiten im Oktavformat, mit einer Kupfertafel.

Der Verfasser hat die Hauptphasen der merkwürdiegen Sonnenfinsternis dieses Jahres für den Nürnberger Horizont, wo ihre Größe 10 Zoll 38 Minuten betragen wird, mit großer Sorgfalt nach der trigonometrischen Methode des Neunzigsten im Voraus berechnet. Einmal ist er dabei ganz der Anordnung von \textsc{Tob. Mayer} gefolgt, das zweite Mal einem etwas, wenn auch nicht wesentlich veränderten Rechengang. Er hat sämtliche Rechenschritte bis ins kleinste Detail in dieser kleinen Schrift abdrucken lassen. Anfängern und ungeübten Liebhabern, die gern eine ähnliche Rechnung für ihren Wohnort ausführen möchten, kann sie als brauchbare und verständliche Anleitung empfohlen werden.
